\section{Introduction}\label{sec:introduction}

Volume does not capture all restrictions on how one body can be embedded
in another by a symplectic map. Such a map preserves a nondegenerate
skew-symmetric form on pairs of tangent vectors, and hence also volume.
Symplectic capacities measure the additional embedding
restrictions, giving a way to compare symplectic size with ordinary volume.
This thesis investigates that comparison for four-dimensional convex
bodies. We develop and use algorithms, numerical experiments and proofs
to study which shapes have large capacity and which deformations can
increase it. AI agents participated in this research, and the project also
examines what it means to delegate mathematical work to them.

We use the Ekeland--Hofer--Zehnder capacity $c_{\rm EHZ}$. For a smooth
convex body it is the least action of a closed characteristic on the
boundary; the preliminaries define these trajectories and their extension
to polytopes. Capacity scales like area, so in four dimensions its square
can be compared with volume through the \emph{systolic ratio}
\[
 \operatorname{sys}(K):=
 \frac{c_{\rm EHZ}(K)^2}{2\operatorname{vol}_4(K)}.
\]
A ball has ratio one. Viterbo's conjecture~\cite{Viterbo2000} asserted that
no convex body has a larger ratio. The bound holds for the Gromov width,
which measures the largest symplectic ball that fits inside a body, simply
by volume preservation. Its failure for another capacity therefore shows
that these notions of symplectic size differ even on convex bodies.

Haim--Kislev and Ostrover found precisely such a failure: a product of two
regular pentagons, placed in the complementary $q$- and $p$-planes, has
systolic ratio $(3+\sqrt5)/5>1$~\cite{HaimKislevOstrover2024}.
We call this the HKO body, $K_{\rm HKO}$; its precise position is given in
Section~\ref{sec:hko-local-maximum}. The counterexample does not identify
the largest possible ratio or explain which shapes make it large. Can the
HKO body be improved by a small deformation? What geometric properties
help identify other bodies with large ratios? These are the main questions
pursued here.

The comparison has consequences for both embeddings and convex geometry.
In four dimensions, EHZ capacity coincides with cylindrical capacity,
which measures the smallest symplectic cylinder into which a body embeds
\cite[Corollary~1]{AbbondandoloEdtmairKang2024}. A connection to convex
geometry concerns how small the product of the areas of a centrally
symmetric planar body $L$ and its polar dual $L^\circ$ can be.
Artstein--Avidan, Karasev and Ostrover prove
$c_{\rm EHZ}(L\times L^\circ)=4$~\cite{AKO2014}, so Viterbo's inequality
on these products becomes the symmetric Mahler inequality
$\operatorname{area}(L)\operatorname{area}(L^\circ)\geq8$.
Section~\ref{subsec:preliminaries-capacity-context} explains these links.

\subsection{Computational and empirical approach}

The definition by closed trajectories does not usually give a capacity
value directly. For polytopes, Haim--Kislev's formula~\cite{HK2017}
turns the problem into a finite quadratic program. We give an expanded
derivation of this formula and an implementation for evaluating capacities.
For products of planar polygons, we prove that at most three facet
directions per factor suffice, making the calculation smaller. A second algorithm
uses the passage geometry of Chaidez and Hutchings~\cite{CH2021}; we prove
its correctness under explicit regularity hypotheses in exact real
arithmetic.

These finite methods make it possible to investigate more than a few
explicitly solvable examples. We sample polytopes, compare geometric
measurements with computed systolic ratios, and test searches that use
capacity evaluations to choose further shapes. Such experiments can reveal
patterns worth explaining and test whether a proposed geometric guide is
useful in a search. The finite formula also supports proofs: feasible
choices in the capacity problem give upper bounds whose variation can
control nearby bodies. Numerical exploration and local analysis thus
address different parts of the capacity--volume question.

\subsection{Results and their scope}

The main local result proves that the HKO body cannot be improved by a
sufficiently small deformation among polytopes with exactly ten facets
(Theorem~\ref{thm:hko-ten-facet-local-maximum}). Equality occurs only under
translations, positive dilations and linear symplectic maps, all of which
preserve the ratio. This gives local optimality at a body with corners
whose ratio exceeds the ball's. For comparison, Abbondandolo and Benedetti
prove local systolic maximality for Zoll contact forms, including the
volume--capacity inequality near the ball~\cite{AB2023}. Our theorem
fixes the number of facets and gives no conclusion for perturbations that
add facets, such as the hyperplane cuts studied by Haim--Kislev
\cite{HaimKislev2025Zoll}. Section~\ref{subsec:hko-related-local-results}
gives the detailed comparison.

There are other local maxima in the same ten-facet space. A rational
hexagon--quadrilateral product has ratio one and a non-strict local
maximum, with a two-parameter family of equality cases after removing the
same symmetries (Theorem~\ref{thm:non-hko-ten-facet-local-maximum}). This
is not another counterexample to Viterbo's inequality; it shows that local
maximality in this space is not confined to the HKO body.

The product-family results control larger deformations within narrower
classes. Independently stretching or shearing the two
regular-pentagon factors cannot improve the HKO ratio, even for large
deformations (Theorem~\ref{thm:pentagon-affine-capacity}); the rotation
family has an explicit capacity formula
(Theorem~\ref{thm:rotated-pentagon-product-formula}). If one factor is a
regular pentagon and the other is any centrally symmetric planar convex
body, we obtain a sharp upper bound strictly below one
(Theorem~\ref{thm:pentagon-symmetric-partner}), so symmetry of the partner
rules out a counterexample in this class.

Even with the planar shapes fixed, the positions of their planes matter.
For a product in the $(q_1,p_1)$- and $(q_2,p_2)$-planes, a small rigid
rotation increases capacity at fixed volume
(Theorem~\ref{thm:product-initial-rotation}). Here the symplectic form is
nondegenerate on each factor plane. This differs from the $q$- and
$p$-planes of HKO, on each of which the form vanishes, so the improving
rotation does not contradict HKO's local maximality.

The empirical study asks whether a geometric measurement can help select
bodies with large ratios. We examine the total symplectic area of the
two-dimensional faces, normalized by $\sqrt{\operatorname{vol}_4}$:
it measures boundary geometry using the same form that defines capacity.
In the sampled populations, smaller values are strongly associated in
rank with larger systolic ratios. This gives a guide to selecting bodies,
but its scope is limited. The association weakens among low-area or
high-ratio bodies. Changing polygon factor shapes by area-preserving
linear maps also weakens it, showing that factor shape contributes to the
pattern. Exact rotation formulas for pentagon and triangle products give
opposite orderings of area and ratio. Thus the experiments identify a
geometric signal, while the family calculations rule out treating it as
a universal prescription for increasing capacity.

A separate search comparison tests procedures that use evaluated
capacities to propose improvements. Methods that accounted for several
competing candidates in the capacity calculation attained larger median
best ratios than literal gradient ascent under common fixed budgets, but later
probes still found directions of improvement missed at their endpoints.
High terminal ratios and failed probes therefore did not certify local
maxima. None of the random samples or search comparisons reported here
produced a numerical ratio above one
(Section~\ref{sec:black-box-datascience}).

\subsection{AI in the research process}

AI use extended to choosing methods and developing mathematical arguments,
software and experiments. In the local maximality work, an agent developed
the argument and exact-computation code from an incomplete suggestion about
controlling nearby bodies by several upper bounds. This provides a
concrete example of delegating mathematical work without prescribing its
solution. Producing readable exposition and keeping agents directed toward
a completed task proved harder. Section~\ref{sec:use-of-ai} reflects on
this experience and on the human judgment needed to direct and assess the
work. For the HKO result, the proof pairs finite exact checks with a
written argument showing why they imply local maximality. The reflection's
observations and hypotheses concern this project; the separate
AI-use disclosure records provenance and the extent of AI use.

\subsection*{Outline}

Sections~\ref{sec:preliminaries}--\ref{sec:quadratic-program-algorithm-hk2019}
develop the capacity formula, Section~\ref{sec:flow-graph-algorithm-ch2021}
the flow-graph algorithm, and Section~\ref{sec:first-order-perturbations}
the first variation. Section~\ref{sec:hko-local-maximum} proves the local
maximality results without using the flow-graph or data chapters.
Section~\ref{sec:black-box-datascience} presents the empirical study, and
Sections~\ref{sec:rotated-regular-polygons}--\ref{sec:product-rotation}
treat the product families. The remaining sections discuss visualization,
numerics, code and data, AI use, and conclusions. The appendices give
the empirical methods and the exact HKO computation.
